\section{3.3~V converter derivation}\label{sec:3v3-converter}

This section derives the required characteristics of the 3.3~V rail and specifies
component values that satisfy them, built around the TPS63070
(\cref{sec:tps63070}). The converter is referred to by part number rather than
by a designator, being the only device of its kind on the rail.
\Cref{sec:hv-converter,sec:esp32-mcu} already designate R206 and the MCU-side
components; the remaining designators are assigned as follows:

\begin{table}[htbp]
\centering
\caption{3.3~V converter: component designators}\label{tab:3v3-designators}
\begin{tabular}{|l|l|}
\hline
\textbf{Designator} & \textbf{Function} \\
\hline
\des{L}{3V3} & Buck-boost inductor \\
\hline
\des{C}{3V3I1}, \des{C}{3V3I2} & Input capacitance \\
\hline
\des{C}{3V3VAUX} & VAUX bypass capacitor \\
\hline
\des{C}{3V3O3}, \des{C}{3V3O2} & Output bulk capacitance \\
\hline
\des{C}{3V3O1} & Output pin bypass capacitor \\
\hline
\des{R}{3V3D1}, \des{R}{3V3D2} & Output-voltage feedback divider \\
\hline
\des{R}{3V3PG} & Power-good pull-up \\
\hline
\des{R}{3V3EN} & Enable and mode pull-up \\
\hline
\des{R}{3V3HDR} & \des{Q}{3V3HDR} gate series resistor \\
\hline
\des{D}{3V3BUS} & Blocking diode, anode on $V_{\mathrm{BUS}}$ \\
\hline
\des{Q}{3V3HDR} & Header-branch ORing switch \\
\hline
\des{R}{BUSBLEED} & Bus-rail discharge path (\cref{sec:3v3-diodeor}) \\
\hline
\des{R}{HDRBLEED} & Header supply-pin discharge path (\cref{sec:3v3-diodeor}) \\
\hline
\end{tabular}
\end{table}

\begin{figure}[htbp]
\centering
\includegraphics[width=\textwidth]{3v3}
\caption{3.3~V converter: reference schematic}\label{fig:3v3}
\end{figure}

\Cref{fig:3v3} shows the connections. \des{R}{BUSBLEED} stands for a
requirement on the bus rail; \cref{sec:3v3-diodeor} derives it and
\cref{sec:pd-bleed} settles how it is met.

The rail's topology comprises: the converter with its inductor, input
capacitance, VAUX bypass and output capacitance, the feedback divider
\des{R}{3V3D1}/\des{R}{3V3D2} setting the output, the MCU
module's supply pin with its local decoupling (\cref{sec:esp32-mcu}), and two
resistive sources feeding other nets, R206 (the flyback's power-good pull-up,
\cref{sec:hv-house-pgood}) and \des{R}{3V3PG}, which terminates this
converter's own power-good output and doubles as the MCU's reset pull-up
(\cref{sec:esp32-reset}). Two things distinguish this rail from the 5~V rail of
\cref{sec:5v-converter}: the load is large and transient-heavy due to Wi-Fi
transmit bursts, and the input is an OR of $V_{\mathrm{BUS}}$ and an externally
supplied header voltage that can sit either above or below the output, which is
what requires a buck-boost rather than a buck.

\subsection{Load budget}\label{sec:3v3-load}

The dominant load is the MCU module. From \cref{sec:esp32}: the worst published
sustained condition is Wi-Fi transmit at \SI{355}{\milli\ampere} (802.11b,
\SI{20.5}{dBm}, rated at 100\,\% duty cycle), and Espressif's own
regulator-sizing requirement is an external supply capable of
$\geq\SI{0.5}{\ampere}$. Typical operation sits far lower: modem-sleep spans
\SIrange{13.2}{91.7}{\milli\ampere} depending on CPU state and clock, with a
\SI{10}{\milli\ampere} adder during flash access; light- and deep-sleep are in
the \si{\micro\ampere} range. The remaining loads:

\begin{itemize}
\item \emph{BCD drive lines.} Sixteen GPIOs drive the four K155ID1s' BCD inputs. A TTL
input driven high draws $I_{IH} \leq \SI{40}{\micro\ampere}$ (A input) or
\SI{80}{\micro\ampere} (B/C/D inputs) from the driving pin (\cref{sec:esp32-bcd}); a
line driven low \emph{sinks} current sourced from the drivers' 5~V rail and adds nothing
to VDD33. Worst case, all sixteen lines high:
\begin{align*}
  I_{\mathrm{BCD}} &= 4 \times (40 + 3\times80)\,\si{\micro\ampere}
    = \SI{1.12}{\milli\ampere},
\end{align*}
0.3\,\% of the Wi-Fi figure: negligible.
\item \emph{Pull-ups.} R206 (the flyback's power-good) and \des{R}{3V3PG} (this
converter's, \cref{sec:3v3-house}) each contribute at most
\SI{33}{\micro\ampere}, and only during power-not-good windows.
\Cref{sec:esp32-pgood,sec:esp32-reset} quantify them.
\item \emph{Serial buffers.} The two logic buffers of \cref{sec:esp32-uart} draw
\SI{4}{\micro\ampere} each at most, plus their idle-state pull-ups at
\SI{330}{\micro\ampere} each while a line is held low.
\end{itemize}

The design totals are therefore
\begin{align*}
  I_{\mathrm{load,max}} &= 355 + 1.12 + 0.03 + 0.7
    \approx \SI{357}{\milli\ampere} \approx \SI{0.36}{\ampere},\\
  P_{\mathrm{out,max}} &= 3.3 \times 0.357 = \SI{1.18}{\watt},
\end{align*}
the sustained worst case, with Espressif's $\geq\SI{0.5}{\ampere}$ capability requirement adopted as the
converter-sizing floor (margin checks below are run at both \SI{0.36}{\ampere} and
\SI{0.5}{\ampere}). For the header-powered scenario of
\cref{sec:3v3-diodeor} (an MCU in download/console mode, radio off), the applicable
draw is the modem-sleep band plus flash access: roughly
\SIrange{30}{100}{\milli\ampere}, with \SI{50}{\milli\ampere} carried as the nominal
estimate. If the module SKU eventually selected has PSRAM, active-mode draw may exceed
the tabulated \SI{355}{\milli\ampere} by an amount Espressif does not quantify
(\cref{sec:esp32}); the headroom between the \SI{0.5}{\ampere} requirement
and the converter's capability absorbs any plausible adder, and the published figures
remain the design point.

\subsection{Operating point}\label{sec:3v3-input}

The converter's input is the OR node of \cref{sec:3v3-diodeor}, which spans
disjoint bands depending on which source is present:
\begin{description}
  \item[Bus-powered, full operation] a \SI{9}{\volt} or \SI{12}{\volt} contract,
    $V_{\mathrm{BUS}}$ \SIrange{8.5}{13}{\volt} through \des{D}{3V3BUS}, giving
    \SIrange{8.1}{12.8}{\volt} at the node. The converter bucks.
  \item[Bus-powered, limited operation] the \SI{5}{\volt} contract of
    \cref{sec:pd-classes}, $V_{\mathrm{BUS}}$ \SIrange{4.75}{5.5}{\volt}, giving
    \SIrange{4.4}{5.15}{\volt} at the node. The converter still bucks. This is
    the band in which the rest of the board is dark and this rail alone carries
    the design, so it is the one band that is not optional.
  \item[Header-powered] an externally supplied \SI{3.3}{\volt} or \SI{5}{\volt}
    through \des{Q}{3V3HDR}, giving \SI{3.15}{\volt} at worst (\cref{sec:3v3-diodeor} derives
    both figures). At \SI{3.15}{\volt} the converter boosts; at \SI{5}{\volt} it
    bucks.
\end{description}
Every band sits inside the part's \SIrange{2.0}{16}{\volt} running range, and the
node's \SI{12.8}{\volt} maximum is $1.6\times$ inside the \SI{20}{\volt} absolute
maximum on VIN. The header-powered band is the one that must also clear the
\SI{3.0}{\volt} start-up requirement, which \cref{sec:3v3-diodeor} treats.

The duty cycles that follow bound two separate checks:
\begin{align*}
  \text{buck: } D &= \frac{V_{OUT}}{V_{IN}}
    &&\SI{12.8}{\volt}: 0.258; \quad \SI{8.1}{\volt}: 0.407; \quad
      \SI{5}{\volt}: 0.661, \\
  \text{boost: } D &= \frac{V_{OUT}-V_{IN}}{V_{OUT}}
    &&\SI{3.15}{\volt}: 0.045.
\end{align*}

\paragraph{Minimum duty cycle in buck mode is the binding constraint at the top
of the bus range.} The part specifies a \SI{30}{\percent} minimum duty in buck
mode, relaxed to \SI{20}{\percent} when the effective output capacitance is at
least \SI{40}{\micro\farad} and the effective inductance falls between
\SIrange{0.7}{1.8}{\micro\henry} (\cref{sec:tps63070}). At the maximum bus input
the duty is \SI{25.8}{\percent}, below the general figure and above the relaxed
one, so the relaxed conditions are not optional: they become requirements on
\des{L}{3V3} and on the output capacitance, imposed in
\cref{sec:3v3-inductor,sec:3v3-cout} below. The requirement would only bind
harder at a higher input; it clears \SI{20}{\percent} for any node voltage below
$3.3/0.20 = \SI{16.5}{\volt}$, above the part's own input maximum.

\subsection{Inductor}\label{sec:3v3-inductor}

The recommended effective inductance is \SIrange{0.7}{2.8}{\micro\henry} with a
nominal of \SI{1.5}{\micro\henry} (\cref{sec:tps63070}), and the minimum-duty
relaxation of \cref{sec:3v3-input} narrows the effective value to
\SIrange{0.7}{1.8}{\micro\henry}. Against the datasheet's anticipated effective
tolerance of $+20/-30\,\%$, a \SI{1.5}{\micro\henry} nominal part spans
\SIrange{1.05}{1.8}{\micro\henry} effective, which fits that narrowed window
exactly. A larger nominal value would not: \SI{2.2}{\micro\henry} reaches
\SI{2.64}{\micro\henry} at its upper extreme and forfeits the relaxation, and
with it the ability to operate at the top of the bus range. Specification:
$\des{L}{3V3} = \SI{1.5}{\micro\henry}$.

\paragraph{Peak current.} Two corners compete. In boost mode the datasheet's own
peak-current expression applies at the minimum input reached there,
\SI{3.15}{\volt}, and is evaluated at the largest current that corner can be
asked for, the full \SI{0.5}{\ampere} capability point:
\begin{align*}
  I_{PEAK} &= \frac{I_{OUT}}{\eta(1-D)} + \frac{V_{IN} D}{2 f L}
    = \frac{0.5}{0.9\times0.955} + \frac{3.15\times0.045}
      {2\times\SI{2.4}{\mega\hertz}\times\SI{1.5}{\micro\henry}}
    = \SI{0.602}{\ampere}.
\end{align*}
In buck mode the peak is the load plus half the ripple, worst at the maximum
input and at the tolerance extremes that maximize ripple (effective inductance
\SI{1.05}{\micro\henry}, oscillator at its \SI{2.1}{\mega\hertz} minimum):
\begin{align*}
  \Delta I_L &= \frac{V_{OUT}(V_{IN}-V_{OUT})}{V_{IN} L f}
    = \frac{3.3\,(12.8-3.3)}{12.8\times\SI{1.05}{\micro\henry}\times\SI{2.1}{\mega\hertz}}
    = \SI{1.11}{\ampere},
  &
  I_{PEAK} &= 0.5 + \frac{1.11}{2} = \SI{1.06}{\ampere}.
\end{align*}
The buck corner governs. Applying the datasheet's recommended \SI{20}{\percent}
saturation margin gives $I_{\mathrm{sat}} \geq \SI{1.27}{\ampere}$ for normal
operation. Against a fault, the converter's own average input current limit
bounds what it can draw at \SI{4.15}{\ampere} maximum, so a part rated at or
above that figure cannot be saturated by the converter under any condition it
can produce. Specification: $I_{\mathrm{sat}} \geq \SI{4.2}{\ampere}$, DCR
$\leq\SI{50}{\milli\ohm}$.

\paragraph{Right-half-plane zero.} Boost operation introduces a right-half-plane
zero, which the datasheet asks to keep above \SI{400}{\kilo\hertz} so it does not
degrade phase margin at crossover. It falls lowest at the minimum input and
maximum load:
\begin{align*}
  f_{RHPZ} &= \frac{(1-D)^2 V_{OUT}}{2\pi I_{OUT} L}
    = \frac{0.955^2 \times 3.3}
      {2\pi \times 0.5 \times \SI{1.5}{\micro\henry}} = \SI{638}{\kilo\hertz},
\end{align*}
clearing the recommendation by $1.6\times$ at the worst corner and by
$2.2\times$ at the \SI{0.357}{\ampere} sustained load. Both figures follow from
the small boost duty cycle: the header-powered case asks the converter to lift
\SI{3.15}{\volt} to \SI{3.3}{\volt}, barely boosting at all.

\subsection{Output voltage programming}\label{sec:3v3-fb}

The adjustable variant regulates its feedback pin to \SI{800}{\milli\volt}, so
the divider ratio sets the output:
\begin{align*}
  \des{R}{3V3D1} &= \des{R}{3V3D2}\left(\frac{V_{OUT}}{\SI{0.8}{\volt}}-1\right)
    = \des{R}{3V3D2}\times 3.125.
\end{align*}
Two effects bound the divider impedance from opposite sides. The feedback pin
draws up to \SI{100}{\nano\ampere} of leakage (\cref{sec:tps63070}); that
current flows in \des{R}{3V3D1} and offsets the output by $I_{lkg}\des{R}{3V3D1}$ irrespective
of the ratio, which favors small resistors. The divider's own current is
quiescent draw that never reaches the load, which favors large ones. The
datasheet bounds the lower leg at \SI{400}{\kilo\ohm} to keep at least
\SI{2}{\micro\ampere} flowing, and recommends lower values than that for
accuracy.

Taking $\des{R}{3V3D2} = \SI{39}{\kilo\ohm}$ puts \SI{20.5}{\micro\ampere} through the
divider, an order of magnitude above the floor and small against the
converter's own \SIrange{54}{103}{\micro\ampere} quiescent current, while
holding the leakage offset to
$\SI{100}{\nano\ampere}\times \des{R}{3V3D1} = \SI{12}{\milli\volt}$. The exact upper
leg is then \SI{121.9}{\kilo\ohm}; the two standard values bracketing it give
\begin{align*}
  \des{R}{3V3D1} = \SI{120}{\kilo\ohm} &\rightarrow V_{OUT} = \SI{3.262}{\volt}, &
  \des{R}{3V3D1} = \SI{130}{\kilo\ohm} &\rightarrow V_{OUT} = \SI{3.467}{\volt}.
\end{align*}
Only the lower is viable: at \SI{3.467}{\volt} the Power Save Mode offset alone
carries the rail to \SI{3.57}{\volt}, and the MCU's supply window has no
headroom above its recommended maximum (\cref{sec:esp32}). Landing below
\SI{3.3}{\volt} is in any case the favorable direction, since that tolerance is
asymmetric at $+3/-1\,\%$ (\cref{sec:tps63070}). Specification:
$\des{R}{3V3D1} = \SI{120}{\kilo\ohm}$, $\des{R}{3V3D2} = \SI{39}{\kilo\ohm}$, both 1\,\% and from
one resistor series so their temperature coefficients track and the ratio does
not drift; $V_{OUT} = \SI{3.262}{\volt}$ programmed.

\subsection{Output capacitance, transient response, and rail accuracy}\label{sec:3v3-cout}

Three requirements bound the output capacitance, all stated in effective terms
because ceramic capacitance falls with DC bias:
\begin{itemize}[nosep]
  \item \emph{Filter ratio.} The adjustable variant requires effective
    capacitance in \si{\micro\farad} at least ten times the effective inductance
    in \si{\micro\henry} (\cref{sec:tps63070}): at the \SI{1.8}{\micro\henry}
    upper extreme, $\geq\SI{18}{\micro\farad}$.
  \item \emph{Minimum duty relaxation.} \Cref{sec:3v3-input} requires
    $\geq\SI{40}{\micro\farad}$ effective. This is the binding requirement.
  \item \emph{Recommended range.} \SIrange{15}{470}{\micro\farad} total at the
    output pin, which bounds the specification from above.
\end{itemize}
\emph{Assumption:} the output ceramics retain $\geq70\,\%$ of nominal
capacitance at \SI{3.3}{\volt} bias. This is a selection requirement on the
parts, not a consequence of their nominal values: capacitance density and bias
derating trade against each other, so a part carrying the whole bulk value in
one body need not meet it where two smaller ones do. Meeting
\SI{40}{\micro\farad} effective then requires \SI{57}{\micro\farad} nominal.

The datasheet's own filter table pairs a \SI{1.5}{\micro\henry} inductor with
\SI{47}{\micro\farad}, \SI{68}{\micro\farad}, or \SI{100}{\micro\farad} of
nominal output capacitance, and separates bulk capacitance from a small-body
part at the output pin (\cref{sec:tps63070}). Specification:
\des{C}{3V3O3} $= \SI{47}{\micro\farad}$ and \des{C}{3V3O2} $= \SI{22}{\micro\farad}$ as bulk,
\SI{69}{\micro\farad} together and so at the table's \SI{68}{\micro\farad}
entry, plus \des{C}{3V3O1} $= \SI{10}{\micro\farad}$ at the pin.
That is \SI{79}{\micro\farad} nominal and \SI{55}{\micro\farad} effective,
clearing the binding requirement by $1.4\times$. Adding the module-local
\SI{22}{\micro\farad} of \cref{sec:esp32-mcu} brings the node to
\SI{101}{\micro\farad} nominal, still inside the \SI{470}{\micro\farad}
ceiling.

\paragraph{Load-step droop (Wi-Fi transmit burst onset).} The bounding step is
idle-to-transmit, $\Delta I = \SI{0.355}{\ampere}$. The inductor's physical slew
limit is not the constraint it was at \SI{1.1}{\mega\hertz}: at the minimum
bus-powered input,
\begin{align*}
  \frac{dI_L}{dt} &= \frac{V_N - V_{OUT}}{L}
    = \frac{8.1-3.3}{\SI{1.5}{\micro\henry}} = \SI{3.2}{\ampere\per\micro\second},
\end{align*}
so the current reaches the new load in \SI{111}{\nano\second}, costing
\SI{20}{\nano\coulomb} and under \SI{1}{\milli\volt}. The dominant term is
control-loop response, for which the datasheet publishes no bandwidth.
\emph{Assumption:} the loop commands full inductor current within
\SI{5}{\micro\second}, carried over from \cref{sec:5v-converter}'s treatment and
conservative for a converter switching more than twice as fast. Then
\begin{align*}
  \Delta V_{\mathrm{droop}} &\approx
    \frac{\Delta I \cdot t_{\mathrm{resp}}}{C_{\mathrm{out}}} + \Delta I \cdot \mathrm{ESR}
    = \frac{0.355 \times \SI{5}{\micro\second}}{\SI{55}{\micro\farad}}
      + 0.355\times\SI{2.5}{\milli\ohm}
    \approx 32 + 1 \approx \SI{33}{\milli\volt}.
\end{align*}
Power Save Mode helps rather than hurts here: it holds the output about
\SI{1}{\percent} high at light load precisely so a step down into PWM has
headroom to fall into.

\paragraph{Accuracy stack against the load's window.} On the adjustable
variant the part's regulation tolerance is specified at the feedback pin, not
at the output (\cref{sec:tps63070}), so the divider's own ratio tolerance and
the feedback leakage of \cref{sec:3v3-fb} add to it at the output rather than
being contained by it. The rail spends most of its
life in Power Save Mode, whose $+3/-1\,\%$ feedback tolerance is the wider of
the two modes and therefore the one to check. The $\pm1\,\%$ divider resistors
carry a $\pm1.5\,\%$ ratio error to the output. Leakage contributes
$\pm\SI{12}{\milli\volt}$. Output ripple in Power Save Mode is
\SI{50}{\milli\volt} pk-pk at \SI{22}{\micro\farad} effective
(\cref{sec:tps63070}), scaling inversely with capacitance to
$\approx\SI{20}{\milli\volt}$ at our \SI{55}{\micro\farad}. Stacking these with
the transient droop above:
\begin{align*}
  V_{\mathrm{rail}}^{\max} &= 0.824\left(1+\frac{121.2}{38.61}\right)
    + 0.012 + 0.020 = \SI{3.44}{\volt}, \\
  V_{\mathrm{rail}}^{\min} &= 0.792\left(1+\frac{118.8}{39.39}\right)
    - 0.012 - 0.020 - 0.033 = \SI{3.12}{\volt},
\end{align*}
against the MCU's \SIrange{3.0}{3.6}{\volt} window (\cref{sec:esp32}):
\SI{116}{\milli\volt} of margin below and \SI{157}{\milli\volt} above. The upper
side is the one that matters most, because the MCU's absolute maximum equals its
recommended maximum with no headroom beyond it, and it is met with the full
Power Save Mode offset, the worst-case divider ratio, leakage in the adverse
direction, and ripple all included.

\subsection{Input capacitance}\label{sec:3v3-cin}

The part requires at least \SI{4.7}{\micro\farad} effective at its input pin,
with \SI{10}{\micro\farad} nominal recommended (\cref{sec:tps63070}), and the
input capacitor must be large enough that supply dips shorter than the
undervoltage-lockout circuit's response time do not reach the pin. At the
\SI{12.8}{\volt} node maximum a \SI{25}{\volt}-rated ceramic loses roughly half
its nominal capacitance to bias, so meeting \SI{4.7}{\micro\farad} effective at
that corner requires about \SI{10}{\micro\farad} nominal. Specification: \des{C}{3V3I1},
\des{C}{3V3I2} $= 2\times\SI{10}{\micro\farad}$, \SI{25}{\volt} X7R or X5R, giving
\SI{20}{\micro\farad} nominal and roughly \SI{10}{\micro\farad} effective at the
worst bias, \des{C}{3V3I2} being the pin-local part and \des{C}{3V3I1} the bulk. The
\SI{12.8}{\volt} node maximum sets a floor of \SI{16}{\volt} on the rating
under the \SI{20}{\percent} derating used throughout; \SI{25}{\volt} is taken
instead so the bias derating assumed above is achievable.

\des{C}{3V3VAUX} $= \SI{100}{\nano\farad}$ bypasses the VAUX pin, the value the datasheet
requires; that pin supplies internal loads only and must carry no other
connection. It is an output of the device rather than a rail, reaching at most
\SI{7}{\volt} (\cref{sec:tps63070}), so its bypass needs
$\geq\SI{10}{\volt}$ under the same derating; the \SI{16}{\volt} grade common
in small bodies meets that.

The buck-boost's input current is discontinuous in boost mode but continuous in
buck mode, and the bus-powered case is always the latter. Average input current
at the minimum bus-powered node voltage is
\begin{align*}
  I_{IN} &= \frac{P_{\mathrm{out,max}}}{\eta\,V_N^{\min}}
    = \frac{1.18}{0.85 \times 8.1} = \SI{0.17}{\ampere}
  \quad (\SI{0.24}{\ampere}\ \text{at the \SI{0.5}{\ampere} capability point}),
\end{align*}
drawn from $V_{\mathrm{BUS}}$ through \des{D}{3V3BUS}, the largest of the logic-rail loads
on the bus budget. \des{D}{3V3BUS} adds $I_{IN} V_F \approx \SI{96}{\milli\watt}$ of series
loss at that point.

\subsection{Input path and header-powered operation}\label{sec:3v3-diodeor}

The converter's input node is fed from two sources: $V_{\mathrm{BUS}}$ through
\des{D}{3V3BUS}, and the serial header's supply pin (\cref{sec:esp32-uart}) through \des{Q}{3V3HDR}.
Each source must be isolated from the other. Without isolation on the header
branch, a bus-powered board would present up to \SI{12.8}{\volt} to the supply
pin of an attached adapter, destroying it; without isolation on the bus branch,
a header-powered board would drive the entire bus rail and its bulk capacitance.
The node follows whichever source is higher, with the other branch blocking.

\paragraph{The start-up requirement sets the header branch's permissible
drop.} The converter needs \SI{3.0}{\volt} at its input to start
(\cref{sec:tps63070}). \emph{Assumption:} a \SI{3.3}{\volt} adapter presents at
least \SI{3.2}{\volt} at the header under the download-mode load, which is a
requirement on external equipment rather than a property of this design and is
restated as such below. The budget between that figure and the converter's
start-up minimum is \SI{200}{\milli\volt}, and the header branch's series
element consumes part of it:
\begin{align*}
  V_N^{\mathrm{hdr}} &= \SI{3.2}{\volt} - I_{IN}R_{DSon}, &
  I_{IN} &= \frac{V_{OUT}I_{OUT}}{\eta V_N}
    = \frac{3.3\times0.1}{0.9\times3.15} = \SI{0.12}{\ampere}
\end{align*}
at the bounding \SI{100}{\milli\ampere} download-mode load. Reserving
\SI{150}{\milli\volt} of margin above the start-up threshold caps the drop at
\SI{50}{\milli\volt}, so $R_{DSon} \leq \SI{0.42}{\ohm}$.

This requirement excludes a diode. A Schottky drops \SIrange{0.25}{0.35}{\volt}
at this current, which alone exceeds the entire \SI{200}{\milli\volt} budget and
would leave the node at roughly \SI{2.9}{\volt}, below the threshold, so the
converter would not start from a \SI{3.3}{\volt} adapter at all. A series
switch is not so constrained. \des{Q}{3V3HDR} is therefore a P-channel
MOSFET, source on the converter's input node and drain on the header supply
pin, so that its body diode conducts from header to node and blocks the reverse
direction that would back-power an adapter.

The \SI{0.42}{\ohm} figure is a target rather than a limit. It follows from
reserving \SI{150}{\milli\volt} of the budget; the point at which the converter
actually fails to start is the whole of it,
$\SI{0.2}{\volt}/\SI{0.12}{\ampere} = \SI{1.67}{\ohm}$. Exceeding the target
erodes reserved margin, exceeding the limit prevents start-up. The distinction
matters because the on-resistance must be met at the gate drive this node
affords, roughly \SI{3}{\volt}, and not at the \SI{4.5}{\volt} such parts are
conventionally characterized at.

\paragraph{The gate references $V_{\mathrm{BUS}}$ rather than ground.} The
switch must be off whenever the bus is present and on whenever it is not, and
tying the gate to the bus rail through \des{R}{3V3HDR} accomplishes both with no control
circuit, because it makes both terminals of the gate-source junction track the
same source. Bus-powered, the gate sits at $V_{\mathrm{BUS}}$ and the node one
\des{D}{3V3BUS} drop below it:
\begin{align*}
  V_{GS} &= V_{\mathrm{BUS}} - \left(V_{\mathrm{BUS}} - V_{F}\right) = +V_{F}
    \approx +\SI{0.35}{\volt},
\end{align*}
positive on a P-channel and therefore off, and off throughout the bus ramp as
well, since gate and node rise together. Header-powered, $V_{\mathrm{BUS}}$
rests at ground and the node at \SI{3.16}{\volt}, giving
$V_{GS} = -\SI{3.15}{\volt}$ and full enhancement.

The largest gate-source excursion anywhere in the design is therefore
\SI{3.15}{\volt}, and the rating requirement is that figure with the
\SI{20}{\percent} derating, $|V_{GS}| \geq \SI{4}{\volt}$, which every
P-channel MOSFET satisfies. A gate referenced to ground instead would put the
full node maximum across that junction and demand
$|V_{GS}| \geq \SI{16}{\volt}$, a rating that is scarce among the
low-threshold parts the drive requirement selects for, since the thin gate
oxide that lowers the threshold is the same oxide that limits the rating.

\des{R}{3V3HDR} $= \SI{100}{\kilo\ohm}$ sits in series with the gate rather than pulling
it anywhere: it works against the transistor's own input capacitance to form a
low pass, so a fast transient on the bus rail reaches the gate slowly while the
node stays where its input capacitance holds it. Gate leakage across it offsets
the gate by \SI{10}{\milli\volt} at \SI{100}{\nano\ampere}, immaterial against
\SI{3.15}{\volt} of drive. The arrangement admits no gate-to-source resistor:
paralleled with \des{R}{3V3HDR} that would divide the node voltage onto the gate and leave
the transistor barely enhanced in the header-powered case.

\paragraph{The gate reference imposes a requirement on the bus rail.} With the
gate tied only to $V_{\mathrm{BUS}}$, that rail is the switch's sole DC
reference and must rest near ground when no bus supply is attached, rather than
floating on residual charge. It does not do so unaided: header-powered, \des{D}{3V3BUS} is
reverse-biased by the node's \SI{3.16}{\volt} and its reverse current flows
into the bus rail. The rail therefore needs a discharge path whose resistance
is bounded by how far $V_{\mathrm{BUS}}$ may rise before \des{Q}{3V3HDR} loses its drive.
Holding $V_{GS}$ at or beyond the $-\SI{2.5}{\volt}$ at which on-resistance is
specified caps $V_{\mathrm{BUS}}$ at \SI{0.66}{\volt}, so
\begin{align*}
  R_{\mathrm{BUSBLEED}} &\leq \frac{\SI{0.66}{\volt}}{I_R},
\end{align*}
with $I_R$ the reverse current of \des{D}{3V3BUS} at \SI{3}{\volt} and the maximum
ambient. As \SI{0.66}{\volt} is a reserved figure rather than a failure point,
so is this: \des{Q}{3V3HDR} does not actually lose enhancement until $V_{GS}$
reaches its threshold maximum, \SI{1.15}{\volt}, which puts $V_{\mathrm{BUS}}$
at \SI{2.0}{\volt} and relaxes the product threefold.

\paragraph{Which term to fix first.} Written this way the bound is a
relationship between two unchosen quantities, and neither is determined by the
other. What breaks the circle is that the two are not equally expensive. The
resistor's cost is a standing draw on the bus rail, $V_{\mathrm{BUS}}/R$,
against a budget of \SI{0.655}{\ampere}; the diode's cost is a search over
reverse current at temperature, where Schottky leakage climbs by more than an
order of magnitude from \SI{25}{\celsius} to the maximum ambient and where
parts characterized at their full reverse rating say little about \SI{3}{\volt}.
The resistor is therefore fixed first and generously:
\begin{align*}
  R_{\mathrm{BUSBLEED}} &= \SI{10}{\kilo\ohm}, &
  \frac{V_{\mathrm{BUS}}}{R} &= \SI{1.28}{\milli\ampere}, &
  P &= \SI{16.4}{\milli\watt}
\end{align*}
at the \SI{12.8}{\volt} top of the window: \SI{0.2}{\percent} of the bus budget,
and inside an 0402's rating by a factor of four. That converts the bound into a
figure a part can be selected against, $I_R \leq \SI{66}{\micro\ampere}$ at
\SI{3}{\volt} and the maximum ambient, against a limit of
\SI{200}{\micro\ampere} at which start-up from the header actually fails. Any
other continuous load the bus rail carries parallels this one and only helps;
the bound is written for the bleeder alone so that it does not depend on
another subsystem retaining its own divider.

The obligation falls on the bus rail, not on this converter, and it is a
bound on impedance rather than a call for a component. Whether the front end
satisfies it as a by-product of its own detach behavior is a question about the
front end, and \cref{sec:pd-bleed} answers it: the port controller's integrated
discharge is an event-driven function with a detach-timing requirement behind
it, not a continuous impedance, and the requirement raised here is continuous.
\des{R}{BUSBLEED} is therefore a populated part. What has to be verified is the
impedance seen at the rail.

\paragraph{A part meeting these bounds.} The PMF170XP
(\ds{PMF170XP.pdf}{6}), a \SI{20}{\volt} P-channel trench MOSFET in SC-70,
satisfies every figure above. It specifies
$R_{DSon} \leq \SI{300}{\milli\ohm}$ at $V_{GS} = -\SI{2.5}{\volt}$ and
$T_j = \SI{25}{\celsius}$, a guaranteed maximum at less gate drive than this
node provides and at eight times the current it carries. Its normalized
temperature curve gives a factor of 1.07 at the \SI{43}{\celsius} junction that
\SI{6}{\milli\watt} of dissipation produces in a \SI{40}{\celsius} ambient,
so \SI{321}{\milli\ohm}, a \SI{38}{\milli\volt} drop, and a header-powered node
at \SI{3.16}{\volt}: \SI{161}{\milli\volt} above the start-up threshold and
inside the reserved margin. The remaining figures clear by wide factors,
$V_{DS} = \SI{20}{\volt}$ against \SI{16}{\volt} required,
$|V_{GS}| = \SI{12}{\volt}$ against \SI{4}{\volt},
$I_D = \SI{1}{\ampere}$ against \SI{0.12}{\ampere}, and a threshold between
\SI{0.65}{\volt} and \SI{1.15}{\volt} against \SI{3.15}{\volt} of drive.

The bus branch faces no such pressure, since a forward drop there is a small
fraction of \SI{8.5}{\volt} and costs only the efficiency accounted in
\cref{sec:3v3-cin}. \des{D}{3V3BUS} is therefore an ordinary Schottky: it must block
$\geq\SI{6.3}{\volt}$ reverse (the \SI{5}{\volt} a header source can place on
the node, with the \SI{20}{\percent} derating used throughout), carry
\SI{0.24}{\ampere} average forward and the \SI{1}{\ampere} the converter's
start-up limit can draw in repetitive peaks, dissipate \SI{96}{\milli\watt},
and hold its reverse current low enough that the bleed resistance above stays
practical. A SOD-123-class part meets the first four; the fifth is what
selection turns on.

\paragraph{The other branch needs the same treatment.} The requirement
\des{R}{BUSBLEED} answers is not peculiar to the bus rail. Each branch of the OR
presents a node that goes undriven whenever the other source is the one running,
and each has a leakage path charging that node from the live side. On the bus
rail it is \des{D}{3V3BUS}'s reverse current; on the header branch it is
\des{Q}{3V3HDR}'s off-state drain leakage, and what it charges is the header's
supply pin itself. Bus-powered with nothing attached, that pin is driven by
nothing, and \SI{1}{\micro\ampere} into a few picofarads of pin and trace takes
it to the node's own voltage in microseconds.

The consequence is not damage. The stored energy at \SI{12.8}{\volt} on that
capacitance is under a nanojoule, and an adapter's own decoupling absorbs it
on contact. It is that a connector pin labeled as a \SI{3.3}{\volt} supply
sits at the top of the bus window, which is a hazard to whatever is connected
to it next and is not what the silkscreen says. Requirement: hold it below one
silicon forward drop, \SI{0.5}{\volt}, so it cannot inject current into an
attached supply rail through any protection structure it meets. With
\des{Q}{3V3HDR}'s leakage specified at \SI{1}{\micro\ampere} maximum at
$T_j = \SI{25}{\celsius}$ and \SI{10}{\micro\ampere} at \SI{150}{\celsius}
(\ds{PMF170XP.pdf}{4}), and the switch dissipating nothing while off so that
$T_j$ is the ambient, log-linear interpolation gives
\SI{1.3}{\micro\ampere} at \SI{40}{\celsius} and
\begin{align*}
  \des{R}{HDRBLEED} &\leq \frac{\SI{0.5}{\volt}}{\SI{1.3}{\micro\ampere}}
    = \SI{385}{\kilo\ohm}.
\end{align*}
Specification: \des{R}{HDRBLEED} $= \SI{100}{\kilo\ohm}$, holding the pin at
\SI{130}{\milli\volt} with a $3.8\times$ margin on the bound, and still inside
it at the \SI{150}{\celsius} figure the part will never see. Attached, it costs
\SI{50}{\micro\ampere} from a \SI{5}{\volt} adapter and
\SI{32}{\micro\ampere} from a \SI{3.3}{\volt} one, three parts in ten thousand
of the \SI{0.12}{\ampere} the branch carries.

It also supplies a condition the rating below assumes. \des{Q}{3V3HDR}'s
$V_{DS}$ requirement is written against the node's full \SI{12.8}{\volt},
which is the drain-to-source difference only if the drain is at ground.
Undriven it is not, and the transistor stands off nothing.

\paragraph{Bus-powered operation.} The node sits at
\SIrange{8.1}{12.8}{\volt} with \des{Q}{3V3HDR}'s gate one \des{D}{3V3BUS} drop above it, so the
transistor is off and its body diode is reverse-biased by the difference
between the node and whatever an attached adapter presents. No forward path
exists from the rail onto an adapter's supply pin, and the design cannot
back-power an adapter.

\paragraph{Header-powered operation, adapter at \SI{3.3}{\volt}.} With
$V_{\mathrm{BUS}}$ absent, \des{D}{3V3BUS} blocks the node from feeding the bus and \des{Q}{3V3HDR}
conducts. The node sits at $\geq\SI{3.15}{\volt}$, above the \SI{3.0}{\volt}
start-up requirement, and the converter boosts to \SI{3.3}{\volt} with a duty
cycle of \SI{4.5}{\percent}. Because the converter regulates rather than passing
the input through, the rail is at its normal \SI{3.3}{\volt} regardless of where
in its tolerance band the adapter sits.

Input current exceeds output current in boost mode, so the adapter must supply
more than the MCU draws: \SI{0.12}{\ampere} at the bounding
\SI{100}{\milli\ampere} download-mode load. Adapters that derive their
\SI{3.3}{\volt} from a bridge chip's internal regulator are commonly limited
near \SI{50}{\milli\ampere} and will not support this; one with its own
regulator generally will. Once started, the converter sustains operation down to
\SI{2.0}{\volt} of input while its output holds above \SI{3.0}{\volt}, so a
sagging adapter degrades gracefully rather than at the start-up threshold.

\paragraph{Header-powered operation, adapter at \SI{5}{\volt}.} The node sits at
\SI{4.99}{\volt}, the converter bucks at a duty cycle of \SI{66}{\percent}, and
the rail regulates identically. This case is unremarkable: it is neither
required nor excluded.

\paragraph{Resulting requirement on the adapter.} The supply pin must present at
least \SI{3.2}{\volt} and supply at least \SI{0.12}{\ampere}; either the
\SI{3.3}{\volt} or the \SI{5}{\volt} setting satisfies the voltage requirement,
and \cref{sec:esp32-uart} leaves the signaling level unconstrained. Attaching or
detaching the adapter while USB-C power is present is safe in either order,
since the node follows the higher source, which is always $V_{\mathrm{BUS}}$.

\subsection{Enable, soft start, and dissipation}\label{sec:3v3-house}

EN and PS/SYNC tie to the converter's input node through
$\des{R}{3V3EN} = \SI{10}{\kilo\ohm}$. The pin's threshold tolerance
(\SIrange{0.77}{0.83}{\volt}, \cref{sec:tps63070}) spans a wider band than the
gap between the converter's \SI{3.0}{\volt} start-up minimum and the
\SI{3.15}{\volt} the input path guarantees, so a divider-programmed lockout
cannot cleanly separate the two; the part's internal undervoltage lockout
together with the source requirements of \cref{sec:3v3-diodeor} keeps operation
inside the specified window. \des{R}{3V3EN} programs nothing, then, and carries
only the pins' \SI{0.2}{\micro\ampere} input current, two microvolts of drop.
It is a pull-up rather than a short so that either pin can be held low
externally without contending with the input node, which is the only means the
board offers of disabling this converter to exercise the other rails alone.
That override is single-valued: with both pins on one node, pulling it low
takes EN with it, so the reachable states are enabled in Power Save Mode and
off, not forced-PWM. PS/SYNC ties high, selecting Power Save Mode: the
rail idles at modem-sleep currents for nearly all of its operating life, and the
accuracy stack of \cref{sec:3v3-cout} absorbs that mode's wider tolerance with
\SI{157}{\milli\volt} to spare. VSEL ties low, leaving the output at the value
\des{R}{3V3D1}/\des{R}{3V3D2} program, and FB2, which VSEL would switch in to raise that value, ties
to ground. The datasheet permits either treatment, ``leave the pin open or
connect to GND if not used'' (\ds{tps63070.pdf}{4}), and ground is taken because
a defined node is preferable to a floating one at a pin that reaches the
feedback network, and because it removes an unterminated input from the
electrical rule check.

PG is an open-drain output, terminated by $\des{R}{3V3PG} = \SI{100}{\kilo\ohm}$ to the
rail. The pull-up belongs to the converter rather than to whatever reads the
signal, so PG presents a defined logic level at this subsystem's own boundary
and does not depend on a resistor placed elsewhere in the design.

Read as a status bit the signal would be worthless: the only device on the
board that could read it draws its own supply from this rail and so cannot be
running to receive a fault indication about that rail. It instead drives the
MCU's enable node (\cref{sec:esp32-reset}), which puts it to work in exactly
the window where the MCU is off. PG releases above
\SIrange{94.5}{98.5}{\percent} of the output and reasserts below
\SIrange{90}{94.5}{\percent} (\cref{sec:tps63070}), which at the programmed
\SI{3.262}{\volt} is \SIrange{3.08}{3.21}{\volt} rising and
\SIrange{2.94}{3.08}{\volt} falling. The rising band sits above the MCU's
\SI{3.0}{\volt} supply minimum (\cref{sec:esp32}), so the module is released
only into a rail that is already in specification. \des{R}{3V3PG} is the only
pull-up on that node (\cref{sec:esp32-reset}), so the current PG sinks while
asserted is $\SI{3.262}{\volt}/\SI{100}{\kilo\ohm} = \SI{33}{\micro\ampere}$,
a thirtieth of the pin's \SI{1}{\milli\ampere} rating, at a saturation voltage
far below the module's $0.25\,V_{DD}$ reset threshold. Its standing load on the
rail is \SI{33}{\micro\ampere}, drawn only while power is not good.

Soft start is internal and load-dependent, \SI{400}{\micro\second} bucking and
\SI{850}{\micro\second} boosting at the datasheet's characterization points,
which are this rail's two operating cases almost exactly. The output-node
capacitance is \SI{101}{\micro\farad} nominal including the module-local part,
\SI{71}{\micro\farad} effective, so the charging current is
\begin{align*}
  I_{\mathrm{chg}} &= C_{\mathrm{out}}\frac{dV}{dt}
    = \SI{71}{\micro\farad}\times\frac{\SI{3.27}{\volt}}{\SI{400}{\micro\second}}
    = \SI{0.58}{\ampere}
\end{align*}
in the faster bucking case, against the \SI{1}{\ampere} start-up current limit
the part applies while power good is low. The MCU draws little until firmware
runs, so the margin holds; were it exceeded, the limit would extend the ramp
rather than fault.

Converter dissipation at the conservative $\eta=0.85$:
\begin{align*}
  P_{\mathrm{diss}} &= P_{\mathrm{out,max}}\left(\frac{1}{\eta}-1\right)
    = 1.18 \times 0.176 = \SI{0.21}{\watt}, &
  \Delta T &= \theta_{JA}P_{\mathrm{diss}} = 63 \times 0.21 \approx \SI{13}{\kelvin},
\end{align*}
rising to \SI{0.29}{\watt} and $\Delta T \approx \SI{18}{\kelvin}$ at the
\SI{0.5}{\ampere} capability point: $T_J \leq \SI{58}{\celsius}$ at a
\SI{40}{\celsius} ambient, far below the \SI{125}{\celsius} junction limit.
Wi-Fi transmit at 100\,\% duty is already the sustained rating, so the
conclusion requires no duty-cycle relief; adding \des{D}{3V3BUS}'s \SI{96}{\milli\watt}
does not change it.

\subsection{Summary}\label{sec:3v3-summary}

Assumptions carried through this section: the header branch's series element
presents $\leq\SI{0.42}{\ohm}$ and an adapter presents $\geq\SI{3.2}{\volt}$ at
the header; a Schottky on the bus branch drops \SIrange{0.3}{0.4}{\volt} over
the operating current range; output ceramics retain $\geq70\,\%$ at
\SI{3.3}{\volt} bias with effective ESR \SI{2.5}{\milli\ohm}, and input
ceramics roughly $50\,\%$ at \SI{12.8}{\volt}; converter efficiency 0.85;
control-loop response to a load step within \SI{5}{\micro\second}; download-mode
MCU draw $\approx\SI{50}{\milli\ampere}$, bounded at \SI{100}{\milli\ampere}.
Each appears at its point of use above.

{\footnotesize
\begin{longtable}{|p{0.23\textwidth}|p{0.335\textwidth}|p{0.335\textwidth}|}
\caption{3.3~V converter: specification summary}\label{tab:3v3-converter-summary}\\
\hline
\textbf{Parameter} & \textbf{Derived requirement} & \textbf{Specification / result} \\
\hline
\endfirsthead
\caption[]{3.3~V converter: specification summary (continued)}\\
\hline
\textbf{Parameter} & \textbf{Derived requirement} & \textbf{Specification / result} \\
\hline
\endhead
\hline
\endlastfoot
Load current & --- & \SI{355}{\milli\ampere} MCU $+$
  $\leq$\SI{2}{\milli\ampere} housekeeping $\approx$ \SI{0.36}{\ampere}
  sustained; capability floor \SI{0.5}{\ampere} \\
\hline
Topology & regulate from inputs above and below \SI{3.3}{\volt} &
  four-switch buck-boost, automatic mode transition \\
\hline
Rail vs.\ MCU window & VDD33 \SIrange{3.0}{3.6}{\volt} &
  \SIrange{3.12}{3.44}{\volt} incl.\ Power Save tolerance, divider tolerance,
  feedback leakage, ripple, droop \\
\hline
Input window & \SIrange{2.0}{16}{\volt} running; $\geq\SI{3.0}{\volt}$
  start-up & node \SIrange{8.1}{12.8}{\volt} bus / $\geq\SI{3.15}{\volt}$
  header \\
\hline
Duty cycle & $\geq\SI{20}{\percent}$ buck (relaxed bound) &
  \SI{25.8}{\percent} at maximum bus input \\
\hline
Inductance \des{L}{3V3} & effective \SIrange{0.7}{1.8}{\micro\henry} for the relaxed
  duty bound & \SI{1.5}{\micro\henry} nominal
  ($\rightarrow$\SIrange{1.05}{1.8}{\micro\henry} effective) \\
\hline
Inductor $I_{\mathrm{sat}}$, DCR & $\geq\SI{1.27}{\ampere}$ normal;
  $\geq\SI{4.2}{\ampere}$ covers the input current limit; DCR
  $\leq\SI{50}{\milli\ohm}$ & $\geq\SI{4.2}{\ampere}$, shielded;
  part to be selected \\
\hline
Right-half-plane zero & $>\SI{400}{\kilo\hertz}$ &
  \SI{638}{\kilo\hertz} at the worst corner \\
\hline
Feedback divider \des{R}{3V3D1}/\des{R}{3V3D2} & $\des{R}{3V3D2}\leq\SI{400}{\kilo\ohm}$;
  \SI{800}{\milli\volt} reference; leakage offset $\propto \des{R}{3V3D1}$ &
  \SI{120}{\kilo\ohm} / \SI{39}{\kilo\ohm} $\rightarrow$
  \SI{3.262}{\volt}; \SI{12}{\milli\volt} leakage offset \\
\hline
Output capacitance & $\geq\SI{40}{\micro\farad}$ effective;
  $\leq\SI{470}{\micro\farad}$ total & $47+22$ bulk $+\ \SI{10}{\micro\farad}$
  at the pin: \SI{55}{\micro\farad} effective \\
\hline
Transmit-burst droop & $\geq\SI{3.0}{\volt}$ at the MCU through the step &
  $\lesssim\SI{33}{\milli\volt}$ droop \\
\hline
Output ripple & --- & $\approx\SI{20}{\milli\volt}$ pk-pk in Power Save
  Mode \\
\hline
Input capacitance \des{C}{3V3I1}, \des{C}{3V3I2} & $\geq\SI{4.7}{\micro\farad}$ effective;
  $\geq\SI{16}{\volt}$ & $2\times\SI{10}{\micro\farad}$, \SI{25}{\volt} \\
\hline
VAUX bypass \des{C}{3V3VAUX} & \SI{100}{\nano\farad}, no other loading;
  $\geq\SI{10}{\volt}$ & \SI{100}{\nano\farad}, \SI{16}{\volt} \\
\hline
Header ORing switch \des{Q}{3V3HDR} & $\leq\SI{0.42}{\ohm}$ at $\leq\SI{3}{\volt}$ gate
  drive (\SI{1.67}{\ohm} is the hard limit); $V_{DS}\geq\SI{16}{\volt}$;
  $|V_{GS}|\geq\SI{4}{\volt}$ & PMF170XP: \SI{321}{\milli\ohm} at
  \SI{43}{\celsius}, node at \SI{3.16}{\volt} \\
\hline
\des{Q}{3V3HDR} gate reference & off with $V_{\mathrm{BUS}}$ present, on without &
  gate to $V_{\mathrm{BUS}}$ via \des{R}{3V3HDR} $=\SI{100}{\kilo\ohm}$;
  $V_{GS} = +0.35$ / $-\SI{3.15}{\volt}$ \\
\hline
Bus rail discharge path & $R_{\mathrm{BUSBLEED}}\leq\SI{0.66}{\volt}/I_R$ so
  the gate reference holds at ground & \SI{10}{\kilo\ohm};
  \SI{1.28}{\milli\ampere} standing draw; sets $I_R\leq\SI{66}{\micro\ampere}$
  on \des{D}{3V3BUS} \\
\hline
Header pin discharge path & $\leq\SI{0.5}{\volt}$ on an undriven supply pin
  against \des{Q}{3V3HDR}'s leakage & \SI{100}{\kilo\ohm}; \SI{130}{\milli\volt}
  at \SI{40}{\celsius}; \SI{50}{\micro\ampere} when attached \\
\hline
Bus blocking diode \des{D}{3V3BUS} & $\geq\SI{6.3}{\volt}$ reverse;
  \SI{0.24}{\ampere} average, \SI{1}{\ampere} peak; \SI{96}{\milli\watt};
  $I_R\leq\SI{66}{\micro\ampere}$ at \SI{3}{\volt}, maximum ambient &
  SOD-123 Schottky class; part to be selected \\
\hline
Backfeed to adapter supply & none permissible & \des{Q}{3V3HDR} off whenever
  $V_{\mathrm{BUS}}$ present \\
\hline
Adapter requirement & $\geq\SI{3.2}{\volt}$ at $\geq\SI{0.12}{\ampere}$ &
  \SI{3.3}{\volt} or \SI{5}{\volt} setting; signaling level unconstrained \\
\hline
Pin disposition & --- & EN, PS/SYNC to input node via
  $\des{R}{3V3EN}=\SI{10}{\kilo\ohm}$ (Power Save);
  VSEL low; FB2 to ground; PG to the MCU enable node \\
\hline
PG as supply supervisor & release only above the MCU's \SI{3.0}{\volt}
  minimum & \des{R}{3V3PG} $=\SI{100}{\kilo\ohm}$; releases
  \SIrange{3.08}{3.21}{\volt}, reasserts \SIrange{2.94}{3.08}{\volt};
  \SI{359}{\micro\ampere} sink \\
\hline
Soft start / inrush & $\leq\SI{1}{\ampere}$ start-up limit & internal
  \SIrange{400}{850}{\micro\second}; \SI{0.58}{\ampere} charging \\
\hline
Converter dissipation & --- & $\leq\SI{0.29}{\watt}$, $\Delta T \leq
  \SI{18}{\kelvin}$ (VQFN, $\theta_{JA}=63$~\si{\celsius\per\watt}) \\
\hline
PSRAM active-mode adder & unquantified (\cref{sec:esp32}) & open margin item;
  absorbed by the capability headroom \\
\hline
\end{longtable}
}

Two findings distinguish this rail from its 5~V sibling. First, its input can
sit either side of its output, which is what selects a buck-boost topology and
in turn makes every subsequent constraint different: the minimum-duty bound at
the top of the bus range propagates into both the inductance and the output
capacitance, and the right-half-plane zero of boost operation has to be checked
at all. Second, the choice of series element on the header branch is forced
rather than incidental: the \SI{200}{\milli\volt} between an adapter's guaranteed output and
the converter's start-up threshold is smaller than any diode's forward drop, so
a controlled switch is required where the bus branch needs only a Schottky.

\subsection{Component selection}\label{sec:3v3-parts}

The requirements derived above, collected as selection criteria. Values are
those specified in this section and ratings are the derived minima. Body sizes
are requirements only at the converter and at the parts the datasheet places
directly on its pins; elsewhere they are the smallest that commonly satisfies
the electrical requirement.

{\footnotesize
\begin{longtable}{|L{0.10\textwidth}|L{0.14\textwidth}|L{0.11\textwidth}|L{0.33\textwidth}|L{0.12\textwidth}|}
\caption{3.3~V converter: component selection}\label{tab:3v3-parts}\\
\hline
\textbf{Designator} & \textbf{Class} & \textbf{Value} & \textbf{Required ratings and characteristics} & \textbf{Body} \\
\hline
\endfirsthead
\caption[]{3.3~V converter: component selection (continued)}\\
\hline
\textbf{Designator} & \textbf{Class} & \textbf{Value} & \textbf{Required ratings and characteristics} & \textbf{Body} \\
\hline
\endhead
\hline
\endlastfoot
--- & Buck-boost converter & TPS63070 & adjustable variant; \SI{800}{\milli\volt} feedback reference & VQFN-13, $2.5\times3$~\si{\milli\meter} \\
\hline
\des{L}{3V3} & Shielded power inductor & \SI{1.5}{\micro\henry} & $I_{\mathrm{sat}}\geq\SI{4.2}{\ampere}$; DCR $\leq\SI{50}{\milli\ohm}$; effective inductance \SIrange{0.7}{1.8}{\micro\henry} including tolerance and current derating; core rated for \SI{2.4}{\mega\hertz} & $\leq5\times5$~\si{\milli\meter}, $\leq\SI{4.5}{\milli\meter}$ high \\
\hline
\des{C}{3V3I1} & Ceramic, input bulk & \SI{10}{\micro\farad} & \SI{25}{\volt}; X5R or X7R; $\geq50\,\%$ capacitance retention at \SI{12.8}{\volt} bias & 0805 \\
\hline
\des{C}{3V3I2} & Ceramic, at VIN pin & \SI{10}{\micro\farad} & as \des{C}{3V3I1} & 0603 \\
\hline
\des{C}{3V3VAUX} & Ceramic, VAUX bypass & \SI{100}{\nano\farad} & $\geq\SI{16}{\volt}$; X7R; nothing else on the pin & 0402 \\
\hline
\des{C}{3V3O3} & Ceramic, output bulk & \SI{47}{\micro\farad} & $\geq\SI{10}{\volt}$; X5R, X6S, or X7R; $\geq70\,\%$ capacitance retention at \SI{3.3}{\volt} bias; ESR $\leq\SI{10}{\milli\ohm}$ & 0805 or 1210 \\
\hline
\des{C}{3V3O2} & Ceramic, output bulk & \SI{22}{\micro\farad} & as \des{C}{3V3O3} & 0805 \\
\hline
\des{C}{3V3O1} & Ceramic, at VOUT pin & \SI{10}{\micro\farad} & as \des{C}{3V3O3} & 0603 \\
\hline
\des{R}{3V3D1} & Film resistor & \SI{120}{\kilo\ohm} & 1\,\%; one series with \des{R}{3V3D2} so the ratio tracks over temperature & 0402 \\
\hline
\des{R}{3V3D2} & Film resistor & \SI{39}{\kilo\ohm} & as \des{R}{3V3D1} & 0402 \\
\hline
\des{R}{3V3PG} & Film resistor & \SI{100}{\kilo\ohm} & $\pm5\,\%$ sufficient; sets only the PG pull-up current & 0402 \\
\hline
\des{R}{3V3HDR} & Film resistor & \SI{100}{\kilo\ohm} & $\pm5\,\%$ sufficient; series in \des{Q}{3V3HDR}'s gate, no gate-to-source resistor alongside it & 0402 \\
\hline
\des{R}{3V3EN} & Film resistor & \SI{10}{\kilo\ohm} & $\pm5\,\%$ sufficient; carries only pin leakage & 0402 \\
\hline
\des{R}{BUSBLEED} & Film resistor & \SI{10}{\kilo\ohm} & $\pm5\,\%$ sufficient; $\geq\SI{62.5}{\milli\watt}$ against \SI{16.4}{\milli\watt} dissipated; populated, the front end's own discharge being detach-timed rather than continuous (\cref{sec:pd-bleed}) & 0402 \\
\hline
\des{R}{HDRBLEED} & Film resistor & \SI{100}{\kilo\ohm} & $\pm5\,\%$ sufficient; header supply pin to ground; carries \si{\micro\ampere}-level leakage while unattached and \SI{50}{\micro\ampere} while attached & 0402 \\
\hline
\des{D}{3V3BUS} & Schottky rectifier & --- & $V_R\geq\SI{6.3}{\volt}$; $I_{F(AV)}\geq\SI{0.3}{\ampere}$; repetitive peak $\geq\SI{1}{\ampere}$; $V_F\leq\SI{0.4}{\volt}$ at \SI{0.24}{\ampere}; $I_R\leq\SI{66}{\micro\ampere}$ at \SI{3}{\volt} and the maximum ambient, set by \des{R}{BUSBLEED} (\cref{sec:3v3-diodeor}); this is the criterion a part chosen for forward drop alone will fail & SOD-123 or smaller \\
\hline
\des{Q}{3V3HDR} & P-channel MOSFET & PMF170XP & $R_{DSon}\leq\SI{0.42}{\ohm}$ \emph{specified at} $\leq\SI{3}{\volt}$ of gate drive; $V_{DS}\geq\SI{16}{\volt}$; $|V_{GS}|\geq\SI{4}{\volt}$; $I_D\geq\SI{0.2}{\ampere}$ continuous; source on the converter node, drain on the header pin & SC-70 (SOT-323) \\
\hline
\end{longtable}
}
